% ===== PRÉAMBULE CANONIQUE — à coller VERBATIM en tête de CHAQUE .tex =====
\documentclass[11pt,a4paper]{article}
\usepackage[utf8]{inputenc}\usepackage[T1]{fontenc}\usepackage{lmodern}
\usepackage[french]{babel}
\usepackage{amsmath,amssymb,mathtools}\usepackage{icomma}
\usepackage[version=4]{mhchem}          % \ce{...} : équations & formules
\usepackage{siunitx}\sisetup{locale=FR} % \qty{}{}, \si{}, \ang{}
\usepackage{chemfig}                    % \chemfig{...} : structures développées
\usepackage{xcolor}\usepackage{colortbl}
\usepackage{tikz}\usepackage{pgfplots}\pgfplotsset{compat=1.18}
\usepackage[most]{tcolorbox}\usepackage{enumitem}\usepackage{array,booktabs}
\usepackage[a4paper,margin=2.2cm]{geometry}
\usetikzlibrary{arrows.meta,calc,angles,quotes,positioning,patterns,backgrounds,shapes.geometric}
\definecolor{primary}{RGB}{222,49,99}\definecolor{primarydark}{RGB}{150,22,55}
\definecolor{defcol}{RGB}{0,114,178}\definecolor{methcol}{RGB}{52,92,148}
\definecolor{excol}{RGB}{0,135,90}\definecolor{attncol}{RGB}{200,40,55}
\tcbset{boxbase/.style={enhanced,breakable,boxrule=0.9pt,arc=2.5pt,
  left=3mm,right=3mm,top=2.3mm,bottom=2.3mm,fonttitle=\bfseries,coltitle=white}}
% --- boîtes TUTORIEL (noms EXACTS, ne pas renommer) ---
\newtcolorbox{cle}[1][]{boxbase,colback=defcol!6,colframe=defcol,
  title={\textsc{À retenir}\ifx\relax#1\relax\relax\else\textmd{ : \itshape #1}\fi}}
\newtcolorbox{methode}[1][]{boxbase,colback=methcol!7,colframe=methcol,sharp corners,
  title={\textsc{Méthode}\ifx\relax#1\relax\relax\else\textmd{ --- \itshape #1}\fi}}
\newtcolorbox{exemplebox}[1][]{boxbase,colback=excol!5,colframe=excol,
  title={\textsc{Exemple}\ifx\relax#1\relax\relax\else\textmd{ : \itshape #1}\fi}}
\newtcolorbox{piege}[1][]{boxbase,colback=attncol!6,colframe=attncol,
  title={\textsc{Piège}\ifx\relax#1\relax\relax\else\textmd{ : \itshape #1}\fi}}
\newtcolorbox{remarque}[1][]{boxbase,colback=black!4,colframe=black!45,
  title={\textsc{Remarque}\ifx\relax#1\relax\relax\else\textmd{ : \itshape #1}\fi}}
% --- boîtes EXERCICE / CORRIGÉ (noms EXACTS, auto-comptées) ---
\newtcolorbox[auto counter]{exo}[1][]{boxbase,colback=primary!4,colframe=primary,
  title={\textsc{Exercice}~\thetcbcounter\ifx\relax#1\relax\relax\else\textmd{ --- \itshape #1}\fi}}
\newtcolorbox[auto counter]{corrige}[1][]{boxbase,colback=excol!4,colframe=excol,
  title={\textsc{Corrigé}~\thetcbcounter\ifx\relax#1\relax\relax\else\textmd{ --- \itshape #1}\fi}}
\newcommand{\R}{\mathbb{R}}\newcommand{\N}{\mathbb{N}}\newcommand{\Z}{\mathbb{Z}}\newcommand{\ds}{\displaystyle}
% =========================================================================
\begin{document}
\sisetup{per-mode=symbol}
\begin{corrige}[$K_{\mathrm{ps}}$ en fonction de $s$]
On applique $K_{\mathrm{ps}}=p^{p}q^{q}s^{p+q}$ (ou on relit le tableau du tutoriel).
\begin{enumerate}[label=\alph*)]
  \item $K_{\mathrm{ps}} = s^{2}$.
  \item $K_{\mathrm{ps}} = 4\,s^{3}$.
  \item $K_{\mathrm{ps}} = 4\,s^{3}$.
  \item $K_{\mathrm{ps}} = 27\,s^{4}$.
  \item $K_{\mathrm{ps}} = 108\,s^{5}$.
\end{enumerate}
\end{corrige}

\begin{corrige}[calculer $K_{\mathrm{ps}}$ à partir de $s$]
On remplace $s$ dans l'expression du type correspondant.
\begin{enumerate}[label=\alph*)]
  \item $K_{\mathrm{ps}} = s^{2} = (\num{1.3e-5})^{2} = \num{1.7e-10}$.
  \item $K_{\mathrm{ps}} = 4s^{3} = 4\,(\num{1.3e-2})^{3} = 4\times\num{2.2e-6} = \num{8.8e-6}$.
  \item $K_{\mathrm{ps}} = 27s^{4} = 27\,(\num{1.0e-10})^{4} = 27\times\num{1.0e-40} = \num{2.7e-39}$.
\end{enumerate}
\end{corrige}

\begin{corrige}[calculer $s$ à partir de $K_{\mathrm{ps}}$ (type $1{:}1$)]
$s=\sqrt{K_{\mathrm{ps}}}$.
\begin{enumerate}[label=\alph*)]
  \item $s = \sqrt{\num{1.8e-10}} = \qty{1.3e-5}{\mole\per\litre}$.
  \item $s = \sqrt{\num{1.1e-10}} = \qty{1.0e-5}{\mole\per\litre}$.
  \item $s = \sqrt{\num{3.3e-9}} = \qty{5.7e-5}{\mole\per\litre}$.
\end{enumerate}
\end{corrige}

\begin{corrige}[calculer $s$ à partir de $K_{\mathrm{ps}}$ (type $1{:}2$ ou $2{:}1$)]
$s=\sqrt[3]{K_{\mathrm{ps}}/4}$.
\begin{enumerate}[label=\alph*)]
  \item $s = \sqrt[3]{\dfrac{\num{8.0e-6}}{4}} = \sqrt[3]{\num{2.0e-6}} = \qty{1.26e-2}{\mole\per\litre}$.
  \item $s = \sqrt[3]{\dfrac{\num{1.1e-12}}{4}} = \sqrt[3]{\num{2.75e-13}} = \qty{6.5e-5}{\mole\per\litre}$.
\end{enumerate}
\end{corrige}

\begin{corrige}[comparer deux sels de même type]
À \emph{même type}, le sel de plus \textbf{petit} $K_{\mathrm{ps}}$ est le moins soluble.
\begin{enumerate}[label=\alph*)]
  \item \ce{AgBr} ($K_{\mathrm{ps}}=\num{5.4e-13}$) est le moins soluble : son $K_{\mathrm{ps}}$ est
        plus petit que celui de \ce{AgCl}.
  \item \ce{BaSO4} ($K_{\mathrm{ps}}=\num{1.1e-10}$) est le moins soluble : son $K_{\mathrm{ps}}$ est
        bien plus petit que celui de \ce{SrSO4}.
\end{enumerate}
\end{corrige}
\end{document}
